%% ============================================================
%%  Introduction — Nature Machine Intelligence variant
%%  Reframed to foreground a PRINCIPLE governing machine
%%  intelligence at scale, rather than a systems-engineering
%%  optimisation. Technical model paragraphs are retained but
%%  the opening, the central question, and the contributions
%%  are recast around "intelligent systems" and "discovery".
%% ============================================================

%% NMI format: the introduction carries no heading.
\label{sec:introduction}

The intelligence of modern foundation models is increasingly delivered not at
training time but at \emph{inference} time, where a trained model is repeatedly
queried to reason, generate, and act.
As these models scale beyond the fast memory of a single accelerator, the
performance of large-scale inference is governed less by arithmetic and more by
how an intelligent system \emph{coordinates data movement} across a hierarchical
memory system spanning device memory, host memory, and secondary
storage~\cite{ma2026inference,gholami2024memorywall}.
Understanding the limits of this coordination is therefore central to
understanding how machine intelligence behaves---and how it can be deployed---at
scale.

This shift carries pressing economic and environmental stakes.
Once a model is trained, the recurring cost of \emph{serving} it dominates its
lifetime resource footprint, and hardware analyses identify memory capacity,
bandwidth, and interconnect latency---not arithmetic throughput---as the binding
constraints of inference~\cite{zhou2024survey}.
The energy penalty of misaligned coordination is substantial: operating in the
I/O-limited regime reduces throughput-per-Watt by more than $3\times$ relative to
coordination-aligned operation, raising CO$_2$ output from 15.8\,g to 42.5\,g per
$10^6$ tokens at the IEA\,(2024) global carbon
intensity~\cite{iea2024electricity}.
Characterising the \emph{intrinsic limits} of memory coordination is thus a
prerequisite for the viability of inference-driven AI.

Existing systems reduce memory pressure via CPU/NVMe
offloading~\cite{flexgen2023,aminabadi2022deepspeedinference},
GPU--CPU swapping~\cite{huang2020swapadvisor},
KV-cache management~\cite{kwon2023pagedattention},
and iteration-level request scheduling in serving engines~\cite{yu2022orca}.
These techniques deliver substantial gains in some configurations yet show
diminishing returns---or even performance inversion---in others.
Prior work explains this heterogeneity through workload-specific factors
(model size, batch size, hardware topology), without a unifying account that
predicts \emph{when} and \emph{why} coordination fails.
This raises a question that is fundamental to how we reason about deploying
intelligent systems: \emph{Is hierarchical memory coordination a continuously
optimisable problem, or does machine intelligence at scale admit intrinsic
boundaries beyond which additional coordination cannot improve---and can
worsen---end-to-end performance?}

We answer this question with a regime-based view.
Denning's working-set model and thrashing theory~\cite{mattson1970evaluation}
identify capacity saturation as a binary event on a \emph{single-tier} system,
and the Roofline model~\cite{williams2009roofline} expresses performance as
$\min(F \cdot I,\, B_{\mathrm{peak}})$ for a fixed workload on a single tier.
Neither framework addresses a \emph{multi-tier} hierarchy in which computation,
locality, cross-tier transfer, and synchronisation co-exist as distinct,
adjustable latency terms---the setting in which large-scale intelligent
inference actually runs.
Our framework---extending Little's Law~\cite{kleinrock1975queueing} to multi-tier
inference---decomposes end-to-end latency into four co-active terms
($T_{\mathrm{comp}} + T_{\mathrm{mem}} + T_{\mathrm{swap}} + T_{\mathrm{sync}}$)
and characterises how their relative magnitudes shift as two dimensionless
ratios change, determining \emph{which} intervention is effective at any
operating point.
Multi-tier inference systems exhibit three operational regimes governed by the
fast-memory residency ratio $R_C = C_{\mathrm{fast}}/W$ and the
transfer-pressure ratio $R_B = B_{\mathrm{slow}} \cdot \Delta t / D$.
Transitions between regimes are abrupt: modest changes in model size, batch
size, or bandwidth shift the dominant latency contributor discontinuously, in a
phase-like manner.
Regime-agnostic strategies effective in one regime fail predictably outside
analytically identifiable boundaries.

We make four contributions:
\begin{itemize}
  \item \textbf{A principle of machine intelligence at scale.}
    We find, empirically and analytically, that multi-tier inference systems
    exhibit three intrinsic operational regimes---capacity-limited,
    coordination-dominated, and I/O-limited---separated by abrupt, phase-like
    transitions, and that the ranking of optimisation strategies \emph{inverts}
    across regime boundaries. This reframes AI-inference scaling as a problem of
    regime-aware design rather than continuous tuning.
  \item \textbf{Minimal analytical model.}
    A structural lower bound proving regime transitions are inevitable;
    $\theta_C = 0.50$ derived analytically from the majority-eviction condition
    (consistent within $\pm 0.02$ across five platforms);
    $\theta_B$ estimable from three measurements in a 10-minute profiling session.
  \item \textbf{Multi-platform validation.}
    Transitions confirmed on NVIDIA A100, TPU v4, AWS Inferentia2, AMD MI250,
    and Intel Optane-PMem, with boundaries consistent to within $\pm 0.05$.
  \item \textbf{Workload generality.}
    Three-regime behaviour demonstrated across autoregressive LLMs, vision QA,
    retrieval-augmented generation, edge detection, and scientific-computing
    workloads (protein structure prediction, climate downscaling), confirming
    applicability to any intelligent workload whose working set exceeds
    fast-memory capacity.
\end{itemize}

Together, these results reframe hierarchical memory coordination as
\emph{regime-aware design} rather than continuous tuning.
To illustrate the framework's practical accessibility, consider Llama-3\,8B at
batch\,=\,8 on a server with 80\,GB HBM and 16.1\,GB/s PCIe bandwidth:
$W = 41.8$\,GB gives $R_C = 80/41.8 = 1.91$ (fully capacity-compliant,
$R_C \gg \theta_C$).
With nearly all weights resident in HBM, compulsory cross-tier transfer is
small ($D \approx 1.2$\,GB per step); at $\Delta t \approx 120$\,ms,
$R_B = 16.1 \times 0.12 / 1.2 = 1.61$ (well inside the coordination-dominated
regime).
By contrast, constraining HBM to 20\,GB forces $R_C = 0.48$ and
$D \approx 20.9$\,GB of eviction traffic per step, pushing the same workload
deep into the capacity-limited regime (see Supplementary Note~1 for the
full quantitative derivation).
Any practitioner can compute these two ratios from hardware datasheets and
model-card parameter counts---before running a single experiment.

\FloatBarrier
